% Candidate rules: functional requirements and assumptions (§3).  Written by hand; see docs/one-id-survey.md.

\section{Functional requirements and assumptions}

\begin{definition}[Requesting the minimum (C-11)]
  \label{def:c11}
  \uses{def:element}
  \notready
  A relying party must request the minimum disclosable attributes and verify
  them.

  A request is a set of data elements, and 'minimum' needs a purpose to be
  measured against. With one, the guarantee is that a request holds nothing
  its purpose does not need.

  Open point: What is the minimum for each use case: booking, advance
  passenger information, an e-gate, a boarding gate? §3.1 says the relying
  party must request the minimum, while §4.2.1.2 says it should request only
  what it needs, and §4.2.1.3 says the claims of a query must be the minimum.
  Which word holds?

  Source:
  \href{https://standards.atlassian.net/wiki/spaces/OneID/pages/2947416082/Technical+Implementation+Guide}{One
  ID guide} §3.1, item 3.
\end{definition}

\begin{definition}[Sharing in full or selectively (C-12)]
  \label{def:c12}
  \uses{def:c11, def:c07}
  \notready
  The passenger must share credentials securely, either in full or as the
  minimum selectively disclosable attributes.

  Both alternatives satisfy the sentence, so it does not limit what a wallet
  sends.

  Open point: May a wallet send a whole credential when the relying party
  asked for part of it?

  Source:
  \href{https://standards.atlassian.net/wiki/spaces/OneID/pages/2947416082/Technical+Implementation+Guide}{One
  ID guide} §3.1, item 4.
\end{definition}

\begin{definition}[Authenticity and integrity (C-13)]
  \label{def:c13}
  \uses{def:c01}
  What a relying party receives must be checked to be authentic and
  unaltered. §4.2 gives the checks for each profile.

  A response is accepted only when the issuer and device checks of its
  profile pass.

  Source:
  \href{https://standards.atlassian.net/wiki/spaces/OneID/pages/2947416082/Technical+Implementation+Guide}{One
  ID guide} §3.1, item 5.
\end{definition}

\begin{definition}[Security considerations (C-14)]
  \label{def:c14}
  \notready
  Security considerations must be followed for the request and the response
  of a presentation.

  Not a rule on its own. The model takes the security considerations to be
  the checks §4.2 lists.

  Open point: Which security considerations are meant: only those written in
  §4.2, or also those of the referenced specifications?

  Source:
  \href{https://standards.atlassian.net/wiki/spaces/OneID/pages/2947416082/Technical+Implementation+Guide}{One
  ID guide} §3.1, item 6.
\end{definition}

\begin{definition}[A list of trusted issuers (C-15)]
  \label{def:c15}
  \uses{def:c01}
  \notready
  The trust model is out of scope, but a relying party must have a mechanism
  to obtain the list of trusted issuers and intermediate certificate
  authorities.

  The list is an input of every issuer check.

  Open point: What is the mechanism, and what does a relying party do when
  the list cannot be obtained or is stale?

  Source:
  \href{https://standards.atlassian.net/wiki/spaces/OneID/pages/2947416082/Technical+Implementation+Guide}{One
  ID guide} §3.2, item 1.
\end{definition}

\begin{definition}[A list of trusted verifiers (C-16)]
  \label{def:c16}
  \uses{def:c01, def:c08}
  \notready
  Wallets must have access to a list of trusted verifiers or readers for
  authentication.

  If the list is required, a wallet refuses a request from a verifier that is
  not on it. The profiles do not say that.

  Open point: Is verifier authentication required? §3.2 says wallets must
  have the list. §4.2.1.2 says a wallet should be able to check verifier
  certificates, calls reader authentication recommended, and notes that not
  every wallet can process it. §4.2.2 says reader authentication should be
  present in a request. May a wallet answer a request it cannot authenticate?

  Source:
  \href{https://standards.atlassian.net/wiki/spaces/OneID/pages/2947416082/Technical+Implementation+Guide}{One
  ID guide} §3.2, item 2.
\end{definition}

\begin{definition}[Secure issuance (C-17)]
  \label{def:c17}
  \uses{def:c03}
  \notready
  Issuance of credentials to the wallet must be secure.

  An assumption with no rule in the guide behind it.

  Open point: What does 'securely issued' require? See C-03.

  Source:
  \href{https://standards.atlassian.net/wiki/spaces/OneID/pages/2947416082/Technical+Implementation+Guide}{One
  ID guide} §3.2, item 3.
\end{definition}

\begin{definition}[All credential types in the wallet (C-18)]
  \label{def:c18}
  \uses{def:photoid, def:boarding-pass, def:c42}
  \notready
  All credential types required for the passenger's journey must be present
  in the holder's wallet.

  A request for two document types can only be answered by a wallet that
  holds both.

  Open point: Which credential types does a journey require, and must they be
  in one wallet? A passenger with a photoid in one wallet and a boarding pass
  in an airline's app holds both and cannot present them together.

  Source:
  \href{https://standards.atlassian.net/wiki/spaces/OneID/pages/2947416082/Technical+Implementation+Guide}{One
  ID guide} §3.2, item 4.
\end{definition}

\begin{definition}[The holder is a passenger (C-19)]
  \label{def:c19}
  \uses{def:c03}
  \notready
  Whoever holds the wallet must be a passenger.

  Nothing in a credential says who holds the wallet. The model keeps the
  subject of each credential as a fact the issuer knows.

  Open point: Must every credential in a wallet be about its holder? A
  parent's wallet may hold the boarding passes of the children travelling
  with them. If that is allowed, holding a boarding pass does not make the
  holder its passenger.

  Source:
  \href{https://standards.atlassian.net/wiki/spaces/OneID/pages/2947416082/Technical+Implementation+Guide}{One
  ID guide} §3.2, item 5.
\end{definition}

\begin{definition}[Discovering contactless travel (C-20)]
  \label{def:c20}
  An airline may be able to find out which airports offer contactless travel.
  Annex B gives informative examples.

  Optional, and outside the credential and presentation rules.

  Source:
  \href{https://standards.atlassian.net/wiki/spaces/OneID/pages/2947416082/Technical+Implementation+Guide}{One
  ID guide} §3.2, item 6.
\end{definition}

\begin{definition}[Directing the passenger to the relying party (C-21)]
  \label{def:c21}
  \uses{def:c20}
  \notready
  Airlines must direct or prompt the passenger to the relying party's web
  page during the journey, if it is available.

  Outside the credential and presentation rules.

  Open point: What makes the page 'available', and what must an airline do
  when several relying parties serve one journey?

  Source:
  \href{https://standards.atlassian.net/wiki/spaces/OneID/pages/2947416082/Technical+Implementation+Guide}{One
  ID guide} §3.2, item 7.
\end{definition}

\begin{definition}[Authentication and consent before every presentation (C-22)]
  \label{def:c22}
  \uses{def:c07}
  \notready
  Every presentation of a credential must happen only after holder
  authentication and consent.

  A response exists only if the holder authenticated and consented to exactly
  the data elements in it.

  Open point: Is holder authentication required before every presentation?
  §3.2 says must. §4.2.1.3 says the response should be released after
  authentication, and §4.2.2 requires consent and does not mention
  authentication.

  Source:
  \href{https://standards.atlassian.net/wiki/spaces/OneID/pages/2947416082/Technical+Implementation+Guide}{One
  ID guide} §3.2, item 8.
\end{definition}

\begin{definition}[Terms and conditions (C-23)]
  \label{def:c23}
  \notready
  Terms and conditions must be explicitly mentioned to the passenger during
  the presentation process.

  Not modelled.

  Open point: Who shows the terms and conditions: the wallet, the relying
  party's page or the airline? Annex B mentions both the airline's and the
  relying party's.

  Source:
  \href{https://standards.atlassian.net/wiki/spaces/OneID/pages/2947416082/Technical+Implementation+Guide}{One
  ID guide} §3.2, item 9.
\end{definition}

\begin{definition}[Keys in a secure area (C-24)]
  \label{def:c24}
  \uses{def:c03}
  \notready
  Keys used for cryptography must be kept in a secure area.

  The model takes a device key to be usable only on the device it was issued
  to.

  Open point: Whose keys, and what counts as a secure area: hardware-backed
  storage on the phone, or any protected store? Does it apply to a verifier's
  request-signing and decryption keys?

  Source:
  \href{https://standards.atlassian.net/wiki/spaces/OneID/pages/2947416082/Technical+Implementation+Guide}{One
  ID guide} §3.2, item 10.
\end{definition}
